\section{Potenciometría}

\resumebox{definicion}{Fuerza magnética del protón}{\footnotesize
    \textbf{Potenciometría:} Técnica que mide la diferencia de potencial entre dos electrodos para determinar la concentración de un ion. Se utiliza comúnmente para medir pH con un electrodo de vidrio.
}

\subsection{Fundamento}
El potencial de un electrodo indicador (ej. electrodo de vidrio para pH) es proporcional al logaritmo de la actividad del ion de interés, según la ecuación de Nernst.

\subsection{Aplicaciones}
\begin{itemize}
    \item Medición de pH en alimentos, agua y suelos.
    \item Determinación de otros iones con electrodos selectivos (ej. $Na^+$, $K^+$).
\end{itemize}

\section{Refractometría}

\resumebox{definicion}{Fuerza magnética del protón}{\footnotesize
    \textbf{Refractometría:} Técnica que mide el índice de refracción de una sustancia, que es una propiedad física relacionada con su concentración y composición.
}

\subsection{Fundamento}
El índice de refracción ($n$) es la relación entre la velocidad de la luz en el vacío y la velocidad de la luz en la sustancia.

\subsection{Aplicaciones}
\begin{itemize}
    \item Determinación de la concentración de sólidos solubles (°Brix) en jugos, mermeladas y otros alimentos.
    \item Control de calidad de aceites, jarabes y productos lácteos.
\end{itemize}

\section{Polarimetría}

\resumebox{definicion}{Fuerza magnética del protón}{\footnotesize
    \textbf{Polarimetría:} Técnica que mide la rotación del plano de luz polarizada por una sustancia ópticamente activa, permitiendo determinar su concentración.
}

\subsection{Fundamento}
Las moléculas quirales rotan el plano de luz polarizada. El ángulo de rotación es proporcional a la concentración del compuesto activo.

\subsection{Aplicaciones}
\begin{itemize}
    \item Determinación de la concentración de azúcares (ej. sacarosa, glucosa) en alimentos.
    \item Determinación de la pureza de compuestos quirales en la industria farmacéutica.
\end{itemize}

\resumebox{ejercicio}{Fuerza magnética del protón}{\footnotesize
    \textbf{Selección de técnica:} ¿Qué técnica (potenciometría, refractometría o polarimetría) usarías para: a) determinar el contenido de azúcar en un jarabe, b) medir el pH de una solución de jugo de limón, y c) verificar la pureza de una muestra de vitamina C?
}